% Bài: L12C1 Bài 2. Nội năng. Định luật I của nhiệt động lực học · Khái niệm nội năng và các yếu tố ảnh hưởng đến nội năng
% ID: L12C1B2-01-DS
% Mức: NB
% ID: L12C1B2-01-DS


% ID: L12C1B2-06-DS
% Mức: NB
% ID: L12C1B2-06-DS
% ID: L12C1B2-06-DS
% Mức: NB
% ID: L12C1B3-80-DS
% ID: L12C1B3-80-DS
% Mức: NB
% ID: L12C1B2-06-DS
%=== Câu 15 ===%
\dangbt{Khái niệm nội năng và các yếu tố ảnh hưởng đến nội năng}
\begin{ex}
% ID: L12C1B2-06-DS
% Mức: NB
Xét các phát biểu sau về dạng “Khái niệm nội năng và các yếu tố ảnh hưởng đến nội năng”
\choiceTF
{\True Nội năng là tổng động năng và thế năng tương tác của các phân tử cấu tạo nên vật.}
{Nội năng của vật bằng cơ năng chuyển động của vật.}
{\True Khi giải dạng “Khái niệm nội năng và các yếu tố ảnh hưởng đến nội năng”, cần xác định đúng trạng thái đầu, trạng thái cuối và quy ước dấu hoặc điều kiện áp dụng.}
{Có thể bỏ qua mọi dữ kiện về khối lượng, nhiệt độ và hiệu suất trong mọi bài toán thuộc dạng này.}
\loigiai{
\begin{itemchoice}
\item (Đúng) Nội năng là tổng động năng và thế năng tương tác của các phân tử cấu tạo nên vật. (Vì): Nội năng là tổng động năng và thế năng tương tác của các phân tử cấu tạo nên vật. b) S: Nội năng của vật bằng cơ năng chuyển động của vật. c) Đ: Khi giải dạng “Khái niệm nội năng và các yếu tố ảnh hưởng đến nội năng”, cần xác định đúng trạng thái đầu, trạng thái cuối và quy ước dấu hoặc điều kiện áp dụng. d) S: Có thể bỏ qua mọi dữ kiện về khối lượng, nhiệt độ và hiệu suất trong mọi bài toán thuộc dạng này. Kết luận dựa trên định nghĩa, điều kiện áp dụng và quy ước trong SGK KNTT
\item (Sai) Nội năng của vật bằng cơ năng chuyển động của vật.
\item (Đúng) Khi giải dạng “Khái niệm nội năng và các yếu tố ảnh hưởng đến nội năng”, cần xác định đúng trạng thái đầu, trạng thái cuối và quy ước dấu hoặc điều kiện áp dụng.
\item (Sai) Có thể bỏ qua mọi dữ kiện về khối lượng, nhiệt độ và hiệu suất trong mọi bài toán thuộc dạng này.
\end{itemchoice}
}
\end{ex}

% ID: L12C1B2-07-DS
% Mức: TH
% ID: L12C1B2-07-DS
% ID: L12C1B2-07-DS
% Mức: TH
% ID: L12C1B3-81-DS
% ID: L12C1B3-81-DS
% Mức: TH
% ID: L12C1B2-07-DS
%=== Câu 17 ===%
\begin{ex}
% ID: L12C1B2-07-DS
% Mức: TH
Xét các phát biểu sau về dạng “Khái niệm nội năng và các yếu tố ảnh hưởng đến nội năng”
\choiceTF
{Nội năng của vật bằng cơ năng chuyển động của vật.}
{\True Khi giải dạng “Khái niệm nội năng và các yếu tố ảnh hưởng đến nội năng”, cần xác định đúng trạng thái đầu, trạng thái cuối và quy ước dấu hoặc điều kiện áp dụng.}
{Có thể bỏ qua mọi dữ kiện về khối lượng, nhiệt độ và hiệu suất trong mọi bài toán thuộc dạng này.}
{\True Nội năng là tổng động năng và thế năng tương tác của các phân tử cấu tạo nên vật.}
\loigiai{
\begin{itemchoice}
\item (Sai) Nội năng của vật bằng cơ năng chuyển động của vật. (Vì): Nội năng của vật bằng cơ năng chuyển động của vật. b) Đ: Khi giải dạng “Khái niệm nội năng và các yếu tố ảnh hưởng đến nội năng”, cần xác định đúng trạng thái đầu, trạng thái cuối và quy ước dấu hoặc điều kiện áp dụng. c) S: Có thể bỏ qua mọi dữ kiện về khối lượng, nhiệt độ và hiệu suất trong mọi bài toán thuộc dạng này. d) Đ: Nội năng là tổng động năng và thế năng tương tác của các phân tử cấu tạo nên vật. Kết luận dựa trên định nghĩa, điều kiện áp dụng và quy ước trong SGK KNTT
\item (Đúng) Khi giải dạng “Khái niệm nội năng và các yếu tố ảnh hưởng đến nội năng”, cần xác định đúng trạng thái đầu, trạng thái cuối và quy ước dấu hoặc điều kiện áp dụng.
\item (Sai) Có thể bỏ qua mọi dữ kiện về khối lượng, nhiệt độ và hiệu suất trong mọi bài toán thuộc dạng này.
\item (Đúng) Nội năng là tổng động năng và thế năng tương tác của các phân tử cấu tạo nên vật.
\end{itemchoice}
}
\end{ex}

% ID: L12C1B2-07-TL
% Mức: NB
% ID: L12C1B2-07-TL
% ID: L12C1B2-07-TL
% Mức: NB
% ID: L12C1B3-82-TL
% ID: L12C1B3-82-TL
% Mức: NB
% ID: L12C1B2-08-TL
%=== Câu 18 ===%
\begin{bt}
% ID: L12C1B2-08-TL
% Mức: NB
Trình bày cách nhận biết và phương pháp giải dạng “Khái niệm nội năng và các yếu tố ảnh hưởng đến nội năng”. Phân tích nhận định sau và sửa lại nếu sai: “Nội năng của vật bằng cơ năng chuyển động của vật.”
\loigiai{
Trước hết xác định hiện tượng, trạng thái và các đại lượng đã cho. Nêu định nghĩa hoặc hệ thức phù hợp, đổi về đơn vị SI, áp dụng đúng điều kiện và kiểm tra ý nghĩa kết quả. Nhận định cần đối chiếu: Nội năng là tổng động năng và thế năng tương tác của các phân tử cấu tạo nên vật. Vì vậy nhận định đã cho sai; phát biểu đúng là: Nội năng là tổng động năng và thế năng tương tác của các phân tử cấu tạo nên vật.
}
\end{bt}

% ID: L12C1B2-08-DS
% Mức: TH
% ID: L12C1B2-08-DS
% ID: L12C1B2-08-DS
% Mức: TH
% ID: L12C1B3-83-DS
% ID: L12C1B3-83-DS
% Mức: TH
% ID: L12C1B2-09-DS
%=== Câu 20 ===%
\begin{ex}
% ID: L12C1B2-09-DS
% Mức: TH
Xét các phát biểu sau về dạng “Khái niệm nội năng và các yếu tố ảnh hưởng đến nội năng” (bộ số 5).
\choiceTF
{Nội năng của vật bằng cơ năng chuyển động của vật.}
{\True Khi giải dạng “Khái niệm nội năng và các yếu tố ảnh hưởng đến nội năng”, cần xác định đúng trạng thái đầu, trạng thái cuối và quy ước dấu hoặc điều kiện áp dụng.}
{Có thể bỏ qua mọi dữ kiện về khối lượng, nhiệt độ và hiệu suất trong mọi bài toán thuộc dạng này.}
{\True Nội năng là tổng động năng và thế năng tương tác của các phân tử cấu tạo nên vật.}
\loigiai{
\begin{itemchoice}
\item (Sai) Nội năng của vật bằng cơ năng chuyển động của vật. (Vì): Nội năng của vật bằng cơ năng chuyển động của vật. b) Đ: Khi giải dạng “Khái niệm nội năng và các yếu tố ảnh hưởng đến nội năng”, cần xác định đúng trạng thái đầu, trạng thái cuối và quy ước dấu hoặc điều kiện áp dụng. c) S: Có thể bỏ qua mọi dữ kiện về khối lượng, nhiệt độ và hiệu suất trong mọi bài toán thuộc dạng này. d) Đ: Nội năng là tổng động năng và thế năng tương tác của các phân tử cấu tạo nên vật. Kết luận dựa trên định nghĩa, điều kiện áp dụng và quy ước trong SGK KNTT
\item (Đúng) Khi giải dạng “Khái niệm nội năng và các yếu tố ảnh hưởng đến nội năng”, cần xác định đúng trạng thái đầu, trạng thái cuối và quy ước dấu hoặc điều kiện áp dụng.
\item (Sai) Có thể bỏ qua mọi dữ kiện về khối lượng, nhiệt độ và hiệu suất trong mọi bài toán thuộc dạng này.
\item (Đúng) Nội năng là tổng động năng và thế năng tương tác của các phân tử cấu tạo nên vật.
\end{itemchoice}
}
\end{ex}

% ID: L12C1B2-08-TLN
% Mức: NB
% ID: L12C1B2-08-TLN
% ID: L12C1B2-08-TLN
% Mức: NB
% ID: L12C1B3-84-TLN
% ID: L12C1B3-84-TLN
% Mức: NB
% ID: L12C1B2-10-TLN
%=== Câu 21 ===%
\begin{ex}
% ID: L12C1B2-10-TLN
% Mức: NB
Cho bốn nhận định: (1) Nội năng là tổng động năng và thế năng tương tác của các phân tử cấu tạo nên vật. (2) Nội năng của vật bằng cơ năng chuyển động của vật. (3) Cần kiểm tra đơn vị và điều kiện áp dụng khi xử lí dạng “Khái niệm nội năng và các yếu tố ảnh hưởng đến nội năng”. (4) Có thể thay tùy ý nhiệt độ Celsius cho nhiệt độ tuyệt đối trong mọi công thức. Có bao nhiêu nhận định đúng? Trả lời bằng một số nguyên.
\shortans{2}
\loigiai{
Nhận định (1) và (3) đúng; (2) và (4) sai. Vậy có 2 nhận định đúng.
}
\end{ex}

% ID: L12C1B2-09-DS
% Mức: VD
% ID: L12C1B2-09-DS
% ID: L12C1B2-09-DS
% Mức: VD
% ID: L12C1B3-85-DS
% ID: L12C1B3-85-DS
% Mức: VD
% ID: L12C1B2-12-DS
%=== Câu 23 ===%
\begin{ex}
% ID: L12C1B2-12-DS
% Mức: VD
Xét các phát biểu sau về dạng “Khái niệm nội năng và các yếu tố ảnh hưởng đến nội năng”
\choiceTF
{\True Khi giải dạng “Khái niệm nội năng và các yếu tố ảnh hưởng đến nội năng”, cần xác định đúng trạng thái đầu, trạng thái cuối và quy ước dấu hoặc điều kiện áp dụng.}
{Có thể bỏ qua mọi dữ kiện về khối lượng, nhiệt độ và hiệu suất trong mọi bài toán thuộc dạng này.}
{\True Nội năng là tổng động năng và thế năng tương tác của các phân tử cấu tạo nên vật.}
{Nội năng của vật bằng cơ năng chuyển động của vật.}
\loigiai{
\begin{itemchoice}
\item (Đúng) Khi giải dạng “Khái niệm nội năng và các yếu tố ảnh hưởng đến nội năng”, cần xác định đúng trạng thái đầu, trạng thái cuối và quy ước dấu hoặc điều kiện áp dụng. (Vì): Khi giải dạng “Khái niệm nội năng và các yếu tố ảnh hưởng đến nội năng”, cần xác định đúng trạng thái đầu, trạng thái cuối và quy ước dấu hoặc điều kiện áp dụng. b) S: Có thể bỏ qua mọi dữ kiện về khối lượng, nhiệt độ và hiệu suất trong mọi bài toán thuộc dạng này. c) Đ: Nội năng là tổng động năng và thế năng tương tác của các phân tử cấu tạo nên vật. d) S: Nội năng của vật bằng cơ năng chuyển động của vật. Kết luận dựa trên định nghĩa, điều kiện áp dụng và quy ước trong SGK KNTT
\item (Sai) Có thể bỏ qua mọi dữ kiện về khối lượng, nhiệt độ và hiệu suất trong mọi bài toán thuộc dạng này.
\item (Đúng) Nội năng là tổng động năng và thế năng tương tác của các phân tử cấu tạo nên vật.
\item (Sai) Nội năng của vật bằng cơ năng chuyển động của vật.
\end{itemchoice}
}
\end{ex}

% ID: L12C1B2-09-TL
% Mức: TH
% ID: L12C1B2-09-TL
% ID: L12C1B2-09-TL
% Mức: TH
% ID: L12C1B3-86-TL
% ID: L12C1B3-86-TL
% Mức: TH
% ID: L12C1B2-13-TL
%=== Câu 24 ===%
\begin{bt}
% ID: L12C1B2-13-TL
% Mức: TH
Trình bày cách nhận biết và phương pháp giải dạng “Khái niệm nội năng và các yếu tố ảnh hưởng đến nội năng”. Phân tích nhận định sau và sửa lại nếu sai: “Nội năng là tổng động năng và thế năng tương tác của các phân tử cấu tạo nên vật.”
\loigiai{
Trước hết xác định hiện tượng, trạng thái và các đại lượng đã cho. Nêu định nghĩa hoặc hệ thức phù hợp, đổi về đơn vị SI, áp dụng đúng điều kiện và kiểm tra ý nghĩa kết quả. Nhận định cần đối chiếu: Nội năng là tổng động năng và thế năng tương tác của các phân tử cấu tạo nên vật. Vì vậy nhận định đã cho đúng.
}
\end{bt}

% ID: L12C1B2-10-DS
% Mức: VDC
% ID: L12C1B2-10-DS
% ID: L12C1B2-10-DS
% Mức: VDC
% ID: L12C1B3-87-DS
% ID: L12C1B3-87-DS
% Mức: VDC
% ID: L12C1B2-14-DS
%=== Câu 26 ===%
\begin{ex}
% ID: L12C1B2-14-DS
% Mức: VDC
Xét các phát biểu sau về dạng “Khái niệm nội năng và các yếu tố ảnh hưởng đến nội năng” 
\choiceTF
{Có thể bỏ qua mọi dữ kiện về khối lượng, nhiệt độ và hiệu suất trong mọi bài toán thuộc dạng này.}
{\True Nội năng là tổng động năng và thế năng tương tác của các phân tử cấu tạo nên vật.}
{Nội năng của vật bằng cơ năng chuyển động của vật.}
{\True Khi giải dạng “Khái niệm nội năng và các yếu tố ảnh hưởng đến nội năng”, cần xác định đúng trạng thái đầu, trạng thái cuối và quy ước dấu hoặc điều kiện áp dụng.}
\loigiai{
\begin{itemchoice}
\item (Sai) Có thể bỏ qua mọi dữ kiện về khối lượng, nhiệt độ và hiệu suất trong mọi bài toán thuộc dạng này. (Vì): Có thể bỏ qua mọi dữ kiện về khối lượng, nhiệt độ và hiệu suất trong mọi bài toán thuộc dạng này. b) Đ: Nội năng là tổng động năng và thế năng tương tác của các phân tử cấu tạo nên vật. c) S: Nội năng của vật bằng cơ năng chuyển động của vật. d) Đ: Khi giải dạng “Khái niệm nội năng và các yếu tố ảnh hưởng đến nội năng”, cần xác định đúng trạng thái đầu, trạng thái cuối và quy ước dấu hoặc điều kiện áp dụng. Kết luận dựa trên định nghĩa, điều kiện áp dụng và quy ước trong SGK KNTT
\item (Đúng) Nội năng là tổng động năng và thế năng tương tác của các phân tử cấu tạo nên vật.
\item (Sai) Nội năng của vật bằng cơ năng chuyển động của vật.
\item (Đúng) Khi giải dạng “Khái niệm nội năng và các yếu tố ảnh hưởng đến nội năng”, cần xác định đúng trạng thái đầu, trạng thái cuối và quy ước dấu hoặc điều kiện áp dụng.
\end{itemchoice}
}
\end{ex}

% ID: L12C1B2-11-TN
% Mức: NB
% ID: L12C1B2-11-TN
% ID: L12C1B2-11-TN
% Mức: NB
% ID: L12C1B3-88-TN
% ID: L12C1B3-88-TN
% Mức: NB
% ID: L12C1B2-15-TN
%=== Câu 29 ===%
\begin{ex}
% ID: L12C1B2-15-TN
% Mức: NB
Câu 1 thuộc dạng “Khái niệm nội năng và các yếu tố ảnh hưởng đến nội năng”. Phát biểu nào sau đây đúng?
\choice
{Nội năng của vật bằng cơ năng chuyển động của vật.}
{Nhiệt lượng và nhiệt độ là hai đại lượng hoàn toàn giống nhau.}
{Mọi quá trình nhiệt học đều không phụ thuộc điều kiện thực hiện.}
{\True Nội năng là tổng động năng và thế năng tương tác của các phân tử cấu tạo nên vật.}
\loigiai{
Phát biểu đúng là: Nội năng là tổng động năng và thế năng tương tác của các phân tử cấu tạo nên vật. Các phương án còn lại trái với mô hình, định nghĩa hoặc hệ thức nhiệt học tương ứng.
}
\end{ex}

% ID: L12C1B2-12-TN
% Mức: NB
% ID: L12C1B2-12-TN
% ID: L12C1B2-12-TN
% Mức: NB
% ID: L12C1B3-89-TN
% ID: L12C1B3-89-TN
% Mức: NB
% ID: L12C1B2-16-TN
%=== Câu 31 ===%
\begin{ex}
% ID: L12C1B2-16-TN
% Mức: NB
Câu 5 thuộc dạng “Khái niệm nội năng và các yếu tố ảnh hưởng đến nội năng”. Phát biểu nào sau đây đúng?
\choice
{Nội năng của vật bằng cơ năng chuyển động của vật.}
{Nhiệt lượng và nhiệt độ là hai đại lượng hoàn toàn giống nhau.}
{Mọi quá trình nhiệt học đều không phụ thuộc điều kiện thực hiện.}
{\True Nội năng là tổng động năng và thế năng tương tác của các phân tử cấu tạo nên vật.}
\loigiai{
Phát biểu đúng là: Nội năng là tổng động năng và thế năng tương tác của các phân tử cấu tạo nên vật. Các phương án còn lại trái với mô hình, định nghĩa hoặc hệ thức nhiệt học tương ứng.
}
\end{ex}

% ID: L12C1B2-13-TN
% Mức: NB
% ID: L12C1B2-13-TN
% ID: L12C1B2-13-TN
% Mức: NB
% ID: L12C1B3-90-TN
% ID: L12C1B3-90-TN
% Mức: NB
% ID: L12C1B2-17-TN
%=== Câu 34 ===%
\begin{ex}
% ID: L12C1B2-17-TN
% Mức: NB
Câu 9 thuộc dạng “Khái niệm nội năng và các yếu tố ảnh hưởng đến nội năng”. Phát biểu nào sau đây đúng?
\choice
{Nội năng của vật bằng cơ năng chuyển động của vật.}
{Nhiệt lượng và nhiệt độ là hai đại lượng hoàn toàn giống nhau.}
{Mọi quá trình nhiệt học đều không phụ thuộc điều kiện thực hiện.}
{\True Nội năng là tổng động năng và thế năng tương tác của các phân tử cấu tạo nên vật.}
\loigiai{
Phát biểu đúng là: Nội năng là tổng động năng và thế năng tương tác của các phân tử cấu tạo nên vật. Các phương án còn lại trái với mô hình, định nghĩa hoặc hệ thức nhiệt học tương ứng.
}
\end{ex}

% ID: L12C1B2-14-TN
% Mức: TH
% ID: L12C1B2-14-TN
% ID: L12C1B2-14-TN
% Mức: TH
% ID: L12C1B3-91-TN
% ID: L12C1B3-91-TN
% Mức: TH
% ID: L12C1B2-18-TN
%=== Câu 37 ===%
\begin{ex}
% ID: L12C1B2-18-TN
% Mức: TH
Câu 2 thuộc dạng “Khái niệm nội năng và các yếu tố ảnh hưởng đến nội năng”. Phát biểu nào sau đây đúng?
\choice
{Nhiệt lượng và nhiệt độ là hai đại lượng hoàn toàn giống nhau.}
{Mọi quá trình nhiệt học đều không phụ thuộc điều kiện thực hiện.}
{\True Nội năng là tổng động năng và thế năng tương tác của các phân tử cấu tạo nên vật.}
{Nội năng của vật bằng cơ năng chuyển động của vật.}
\loigiai{
Phát biểu đúng là: Nội năng là tổng động năng và thế năng tương tác của các phân tử cấu tạo nên vật. Các phương án còn lại trái với mô hình, định nghĩa hoặc hệ thức nhiệt học tương ứng.
}
\end{ex}

% ID: L12C1B2-15-TN
% Mức: TH
% ID: L12C1B2-15-TN
% ID: L12C1B2-15-TN
% Mức: TH
% ID: L12C1B3-92-TN
% ID: L12C1B3-92-TN
% Mức: TH
% ID: L12C1B2-19-TN
%=== Câu 40 ===%
\begin{ex}
% ID: L12C1B2-19-TN
% Mức: TH
Câu 6 thuộc dạng “Khái niệm nội năng và các yếu tố ảnh hưởng đến nội năng”. Phát biểu nào sau đây đúng?
\choice
{Nhiệt lượng và nhiệt độ là hai đại lượng hoàn toàn giống nhau.}
{Mọi quá trình nhiệt học đều không phụ thuộc điều kiện thực hiện.}
{\True Nội năng là tổng động năng và thế năng tương tác của các phân tử cấu tạo nên vật.}
{Nội năng của vật bằng cơ năng chuyển động của vật.}
\loigiai{
Phát biểu đúng là: Nội năng là tổng động năng và thế năng tương tác của các phân tử cấu tạo nên vật. Các phương án còn lại trái với mô hình, định nghĩa hoặc hệ thức nhiệt học tương ứng.
}
\end{ex}

% ID: L12C1B2-16-TN
% Mức: TH
% ID: L12C1B2-16-TN
% ID: L12C1B2-16-TN
% Mức: TH
% ID: L12C1B3-93-TN
% ID: L12C1B3-93-TN
% Mức: TH
% ID: L12C1B2-20-TN
%=== Câu 42 ===%
\begin{ex}
% ID: L12C1B2-20-TN
% Mức: TH
Câu 10 thuộc dạng “Khái niệm nội năng và các yếu tố ảnh hưởng đến nội năng”. Phát biểu nào sau đây đúng?
\choice
{Nhiệt lượng và nhiệt độ là hai đại lượng hoàn toàn giống nhau.}
{Mọi quá trình nhiệt học đều không phụ thuộc điều kiện thực hiện.}
{\True Nội năng là tổng động năng và thế năng tương tác của các phân tử cấu tạo nên vật.}
{Nội năng của vật bằng cơ năng chuyển động của vật.}
\loigiai{
Phát biểu đúng là: Nội năng là tổng động năng và thế năng tương tác của các phân tử cấu tạo nên vật. Các phương án còn lại trái với mô hình, định nghĩa hoặc hệ thức nhiệt học tương ứng.
}
\end{ex}

% ID: L12C1B2-17-TN
% Mức: VD
% ID: L12C1B2-17-TN
% ID: L12C1B2-17-TN
% Mức: VD
% ID: L12C1B3-94-TN
% ID: L12C1B3-94-TN
% Mức: VD
% ID: L12C1B2-74-TN
%=== Câu 44 ===%
\begin{ex}
% ID: L12C1B2-74-TN
% Mức: VD
Câu 3 thuộc dạng “Khái niệm nội năng và các yếu tố ảnh hưởng đến nội năng”. Phát biểu nào sau đây đúng?
\choice
{Mọi quá trình nhiệt học đều không phụ thuộc điều kiện thực hiện.}
{\True Nội năng là tổng động năng và thế năng tương tác của các phân tử cấu tạo nên vật.}
{Nội năng của vật bằng cơ năng chuyển động của vật.}
{Nhiệt lượng và nhiệt độ là hai đại lượng hoàn toàn giống nhau.}
\loigiai{
Phát biểu đúng là: Nội năng là tổng động năng và thế năng tương tác của các phân tử cấu tạo nên vật. Các phương án còn lại trái với mô hình, định nghĩa hoặc hệ thức nhiệt học tương ứng.
}
\end{ex}

% ID: L12C1B2-18-TN
% Mức: VD
% ID: L12C1B2-18-TN
% ID: L12C1B2-18-TN
% Mức: VD
% ID: L12C1B3-95-TN
% ID: L12C1B3-95-TN
% Mức: VD
% ID: L12C1B2-75-TN
%=== Câu 46 ===%
\begin{ex}
% ID: L12C1B2-75-TN
% Mức: VD
Câu 7 thuộc dạng “Khái niệm nội năng và các yếu tố ảnh hưởng đến nội năng”. Phát biểu nào sau đây đúng?
\choice
{Mọi quá trình nhiệt học đều không phụ thuộc điều kiện thực hiện.}
{\True Nội năng là tổng động năng và thế năng tương tác của các phân tử cấu tạo nên vật.}
{Nội năng của vật bằng cơ năng chuyển động của vật.}
{Nhiệt lượng và nhiệt độ là hai đại lượng hoàn toàn giống nhau.}
\loigiai{
Phát biểu đúng là: Nội năng là tổng động năng và thế năng tương tác của các phân tử cấu tạo nên vật. Các phương án còn lại trái với mô hình, định nghĩa hoặc hệ thức nhiệt học tương ứng.
}
\end{ex}

% ID: L12C1B2-19-TN
% Mức: VDC
% ID: L12C1B2-19-TN
% ID: L12C1B2-19-TN
% Mức: VDC
% ID: L12C1B3-96-TN
% ID: L12C1B3-96-TN
% Mức: VDC
% ID: L12C1B2-76-TN
%=== Câu 49 ===%
\begin{ex}
% ID: L12C1B2-76-TN
% Mức: VDC
Câu 4 thuộc dạng “Khái niệm nội năng và các yếu tố ảnh hưởng đến nội năng”. Phát biểu nào sau đây đúng?
\choice
{\True Nội năng là tổng động năng và thế năng tương tác của các phân tử cấu tạo nên vật.}
{Nội năng của vật bằng cơ năng chuyển động của vật.}
{Nhiệt lượng và nhiệt độ là hai đại lượng hoàn toàn giống nhau.}
{Mọi quá trình nhiệt học đều không phụ thuộc điều kiện thực hiện.}
\loigiai{
Phát biểu đúng là: Nội năng là tổng động năng và thế năng tương tác của các phân tử cấu tạo nên vật. Các phương án còn lại trái với mô hình, định nghĩa hoặc hệ thức nhiệt học tương ứng.
}
\end{ex}

% ID: L12C1B2-20-TN
% Mức: VDC
% ID: L12C1B2-20-TN
% ID: L12C1B2-20-TN
% Mức: VDC
% ID: L12C1B3-97-TN
% ID: L12C1B3-97-TN
% Mức: VDC
% ID: L12C1B2-77-TN
%=== Câu 52 ===%
\begin{ex}
% ID: L12C1B2-77-TN
% Mức: VDC
Câu 8 thuộc dạng “Khái niệm nội năng và các yếu tố ảnh hưởng đến nội năng”. Phát biểu nào sau đây đúng?
\choice
{\True Nội năng là tổng động năng và thế năng tương tác của các phân tử cấu tạo nên vật.}
{Nội năng của vật bằng cơ năng chuyển động của vật.}
{Nhiệt lượng và nhiệt độ là hai đại lượng hoàn toàn giống nhau.}
{Mọi quá trình nhiệt học đều không phụ thuộc điều kiện thực hiện.}
\loigiai{
Phát biểu đúng là: Nội năng là tổng động năng và thế năng tương tác của các phân tử cấu tạo nên vật. Các phương án còn lại trái với mô hình, định nghĩa hoặc hệ thức nhiệt học tương ứng.
}
\end{ex}

% Mức: NB
% ID: L12C1B2-79-TN
% ID: L12C1B2-79-TN
% Mức: NB
% ID: L12C1B3-98-TN
% ID: L12C1B3-98-TN
% Mức: NB
% ID: L12C1B2-79-TN
%=== Câu 87 ===%
\begin{ex}
% ID: L12C1B2-79-TN
% Mức: NB
	Phát biểu nào sau đây nói về nội năng là \textbf{không} đúng?
	\choice 
	{Nội năng của một hệ có thể tăng lên hoặc giảm xuống}
	{Nội năng là một dạng năng lượng}
	{Nội năng có thể chuyển hóa thành các dạng năng lượng khác}
	{\True Nội năng là nhiệt lượng}
	\loigiai{
		Nội năng là tổng động năng và thế năng của các phân tử cấu tạo nên vật, nên nội năng là một dạng năng lượng. Nội năng của một hệ có thể tăng hoặc giảm và cũng có thể chuyển hóa thành các dạng năng lượng khác. 
		
		Nhiệt lượng không phải là nội năng. Nhiệt lượng là phần năng lượng truyền từ vật này sang vật khác do có sự chênh lệch nhiệt độ.
		
		Vì vậy phát biểu \textbf{“Nội năng là nhiệt lượng”} là \textbf{không đúng}.
	}
\end{ex}

% Mức: TH
% ID: L12C1B2-80-TN
% ID: L12C1B2-80-TN
% Mức: TH
% ID: L12C1B3-99-TN
% ID: L12C1B3-99-TN
% Mức: TH
% ID: L12C1B2-80-TN
%=== Câu 88 ===%
\begin{ex}
% ID: L12C1B2-80-TN
% Mức: TH
	 Nếu làm tăng nhiệt độ của một hệ mà không làm thay đổi thể tích của nó thì nội năng của hệ sẽ
	 \choice 
	 {\True tăng lên}
	 {không thay đổi}
	 {giảm xuống}
	 {tăng lên sau đó giảm xuống}
	\loigiai{
	Khi nhiệt độ tăng thì động năng
	}
\end{ex}
